
\section{Übersicht verwendeter Hilfsmittel}
\label{sec:hilfsmittel}
In Übereinstimmung mit der Erklärung zur Projektarbeit werden nachfolgend die im Rahmen dieser Arbeit genutzten KI-Werkzeuge sowie die Art und Weise ihrer Verwendung offengelegt.

\subsection*{Eingesetztes KI-Werkzeug}
\begin{itemize}
    \item \textit{Produktname}: Gemini (Google)
    \item \textit{Modellversion}: Gemini 1.5 Flash (Stand: Juli 2026)
\end{itemize}

\subsection*{Art und Umfang der Nutzung}
Das KI-Werkzeug wurde während der Erstellung der Projektarbeit als assistierendes Werkzeug für die Textstrukturierung, die sprachliche Formulierung und die redaktionelle Überarbeitung eingesetzt. Die inhaltlichen, technischen und empirischen Ergebnisse (wie Messdaten, Filterfrequenzen, Modellentscheidungen und Beobachtungen) basieren vollständig auf der Eigenleistung des Autors. Die KI wurde in folgenden Bereichen unterstützend genutzt:
\begin{itemize}
    \item \textit{Gliederung und Strukturierung}: Unterstützung bei der logischen Anordnung der Unterkapitel (z. B. Zusammenführung von Feature-Selektion und Modellwahl in Kapitel 5.3 vor der Firmware-Implementierung).
    \item \textit{Ausformulierung von Stichpunkten}: Überführung von stichpunktartigen, technischen Notizen und mathematischen Zusammenhängen des Autors in einen fließenden, akademischen Text.
    \item \textit{Textoptimierung}: Iterative Verfeinerung von Absätzen (z. B. präzisere Einbindung der physikalischen Einheiten, der Erklärung zum Random Forest Classifier sowie der Erläuterung von Aliasing-Effekten beim Downsampling).
\end{itemize}

\subsection*{Übersicht der genutzten Eingabeparameter (Beispiel-Prompts)}
Die Interaktion mit dem KI-Werkzeug erfolgte über eine fortlaufende, interaktive Chat-Sitzung. Die Prompts wurden vom Autor situativ angepasst. Nachfolgend sind Beispiele für wesentlichen Kern-Prompts aufgeführt, welche die Struktur und Tonalität der generierten Textpassagen maßgeblich gesteuert haben:
\begin{itemize}
    \item Hier ist die überarbeitete und präzisere Fassung für den Punkt Randomisierung: [\dots] Bitte den Punkt trotzdem knapp halten und im Fließtext formulieren, sodass Stichpunkte vermieden werden.
    \item Ich denke, wir sollten das Kapitel 5.4 Modellwahl und Training vorziehen. Vielleicht dieses Kapitel umbennenen in "'Feature-Selektion, Modellwahl und Training"'. Bitte Stichpunkte/Aufzählungen nur da wo sinnvoll und sonst eher Fließtext und mathematische Formeln nutzen.“
    \item Hiernach wünsche ich mir noch eine kurze Erklärung zu Random Forest Classifiers, was sind die, wie funktionieren sie? Sehr kurz und als Fließtext.
    \item Wichtig hierbei ist, dass die komplette mathematische Signalbereinigung durch die Biquad-Filter fundamental bei der vollen nativen Abtastrate von 1000 Hz durchgeführt wird. Bitte bei der Neuformulierung möglichst die Originalstruktur beibehalten und diesen Umstand an einer geeigneten Stelle einfügen.
    \item Bitte fasse die folgenden empirischen Beobachtungen (Systemlatenz von 0,5s, biomechanische Unterschiede der Gesten, Konfusionen zwischen Schere und Papier, Verhalten bei Sensor-Repositionierung) in Kapitel 6 als reinen, wissenschaftlichen Fließtext ohne Stichpunkte zusammen.
\end{itemize}
